\headline{{\textbf{\Large\sffamily Seasonal Care Tips:}}\\
          {\textbf{\large\sffamily Mid\--March to Mid\--April}}}
\label{2025-03-care}
\noindent
As the days continue to lengthen and temperatures gradually rise, mid\--March to mid\--April marks the true beginning of spring. 
Trees that have been dormant over winter are now waking up, with buds swelling and new growth emerging. 
This is an exciting and busy time for bonsai enthusiasts, as there are plenty of essential tasks to ensure a strong and healthy growing season. 
Whether you're preparing trees for a show, developing their structure, or simply keeping them in good health, this guide will help you navigate the early spring period with confidence.

\medskip

\topic{Monitoring Spring Growth}
\noindent
With the increasing daylight and warmer temperatures, your bonsai will be putting out new shoots. 
Deciduous trees like maples, elms, and beeches are leafing out, while evergreens such as pines and junipers are also beginning to show signs of activity. 

Keep an eye on this fresh growth, as it's a great indicator of your tree's health. 
Vigorous growth is a sign that your tree is happy, while weak or delayed budding may suggest issues with root health, watering, or nutrient availability. 
If you notice a tree struggling, now is the time to investigate and adjust your care routine accordingly.

\medskip

\topic{Watering and Feeding}
\noindent
As temperatures rise, your trees will start using more water. 
Be mindful that soil dries out faster now compared to winter, but avoid falling into a rigid watering routine. 
Check the moisture levels regularly---if the top centimetre of soil feels dry, it's time to water. 

This is also the right time to introduce feeding.
Start with a balanced fertiliser, preferably organic, to support healthy leaf and root development.
Deciduous trees benefit from a nitrogen\--rich feed to encourage early growth, while conifers prefer a more balanced mix.
Feed moderately at first, increasing frequency as the season progresses.

\medskip

\topic{Wiring and Pruning}
\noindent
Spring is ideal for fine pruning and shaping.
Deciduous trees can be pruned now to refine their structure, particularly if there are any unwanted shoots disrupting the silhouette.
Pinch back new growth on species like maples to encourage finer ramification.

For evergreens, wiring can be done carefully, but avoid excessive bending, as the branches are more brittle at this time of year.
When applying wire, keep an eye on swelling buds---branches thicken quickly in spring, and wire can bite in faster than expected.

\medskip

\topic{Repotting: The Final Window}
\noindent
If you haven't repotted your bonsai yet, this is your last chance for most species before active growth is in full swing.
Trees such as maples, larches, and hornbeams should be repotted before their new leaves harden.

If you do repot, make sure to keep the tree in a sheltered location for a couple of weeks afterwards to allow the roots to recover.
Mist the foliage regularly, and avoid strong sun or drying winds that could stress the tree during this crucial period.

\medskip

\topic{Pest and Disease Prevention}
\noindent
As the weather warms, pests and fungal issues become more prevalent.
Aphids, scale insects, and spider mites can quickly take hold if left unchecked.
Check the undersides of leaves and along branches for any signs of pests, and remove them as soon as possible.

A light spray of neem oil or insecticidal soap can help keep pests under control.
Also, be mindful of fungal infections, particularly on species prone to mildew or black spot.
Good airflow around your trees, as well as avoiding overhead watering in the evenings, will help prevent these problems.

\medskip

\topic{Sunlight and Positioning}
\noindent
Now that the worst of winter is behind us, you can gradually start moving your bonsai back into their regular outdoor positions.
Deciduous trees, in particular, appreciate more sunlight as they leaf out.

Be careful with late frosts, though!
In London, cold snaps can still occur in early spring.
If frost is forecast, bring vulnerable species under shelter or cover them with horticultural fleece overnight.

\medskip

\topic{Spring Shows and Display}
\noindent
For those preparing trees for spring exhibitions, now is the time for final refinements.
Clean the bark, ensure moss is neatly arranged, and check the overall presentation of the pot and tree.

If you are displaying a bonsai indoors temporarily, remember that they should not remain inside for extended periods.
A few days is fine, but prolonged indoor conditions can weaken them.

\medskip

\topic{Enjoy the Season}
\noindent
Mid\--March to mid\--April is a period of transformation, as bonsai shift from dormancy into a season of vigorous growth.
There's a lot to do, but this is also one of the most rewarding times of the year for bonsai enthusiasts.

Take the time to appreciate the changes happening in your trees.
Whether it's the delicate unfurling of a maple leaf or the subtle deepening of pine needle colour, these small moments are what make bonsai cultivation so special.

Happy bonsai tending!

\closearticle

\columnbreak
